\documentclass[12pt]{article}
\usepackage{report}

\begin{document}

\reporttitle{CS340400 Compiler Design - Homework 3 Bonus Implementation Report}

\section{Overview}

The Bonus extends the HW3 code generator to a 32-bit RISC-V target and adds four SIMD intrinsics from
the Andes P-extension (DSP), so that the generated code can run on an Andes Corvette-F1/T1 board through
the Arduino IDE.

\begin{center}
\begin{tabularx}{0.95\linewidth}{@{}lL@{}}
\toprule
Files & the same \file{scanner.l} and \file{parser.y} as HW3, as the specification requires \\
New features & \code{--rv32} target, the \code{uint32_t} type with unsigned operations, 4 intrinsics \\
RV32 & all 15 HW3 test programs compiled with \code{--rv32} on RV32 Spike --- \textbf{15/15 pass} \\
Bonus tests & 5 programs in \file{bonus/testcases/}, 12 inputs each, with the P-extension emulated ---
             \textbf{5/5 pass} \\
Not yet verified & running on the board; whether the Andes assembler accepts the output \\
Usage & \code{./codegen --rv32 < testcase.c} \\
\bottomrule
\end{tabularx}
\end{center}

\subsection{Requirements}

\begin{itemize}
  \item Use the \emph{same} \file{scanner.l} and \file{parser.y} as HW3 and switch to 32-bit code with a
        flag; RV64 stays the default.
  \item On RV32, memory accesses are 4 bytes wide, with \code{lw} and \code{sw}.
  \item Support four intrinsics that operate on the four bytes of a 32-bit value independently:
\end{itemize}

\begin{center}
\begin{tabularx}{\linewidth}{@{}llL@{}}
\toprule
Intrinsic & Instruction & Per-byte result \\
\midrule
\code{__rv__ukadd8(a, b)} & \code{ukadd8 rd, rs1, rs2} & unsigned saturating add: $\min(a_i + b_i,\ 255)$ \\
\code{__rv__uksub8(a, b)} & \code{uksub8 rd, rs1, rs2} & unsigned saturating subtract: $\max(a_i - b_i,\ 0)$ \\
\code{__rv__cmpeq8(a, b)} & \code{cmpeq8 rd, rs1, rs2} & \code{0xFF} if $a_i = b_i$, else \code{0x00} \\
\code{__rv__ucmplt8(a, b)} & \code{ucmplt8 rd, rs1, rs2} & \code{0xFF} if $a_i < b_i$ (unsigned), else \code{0x00} \\
\bottomrule
\end{tabularx}
\end{center}

A test program defines \code{uint32_t ext_dsp_codegen(uint32_t a, uint32_t b)}; the TA's sketch
\file{HW3_Bonus.ino} reads two numbers from the serial port, calls it, and compares the result with
the board's built-in intrinsics. The TA's example:

\lstinputlisting[title={\codetitle{bonus/testcases/0.c}}]{../testcases/0.c}

\section{Selecting the Target}

HW3 already routes every size-dependent decision through five variables, so the RV32 target only has
to change them:

\snippet[language=yacc]{../../parser.y}{static int W = 8;}{+5}
\snippet[language=yacc]{../../parser.y}{static void use_rv32}{@end}
\snippet[language=yacc]{../../parser.y}{int main(}{@end}

\code{--rv32} selects the 32-bit target. The environment variable \code{CODEGEN_TARGET=rv32} does the
same, so a grading script that runs plain \code{./codegen < t.c} can still be pointed at RV32 without
changing its command line. \code{--pext-emulate} is a test-only option described in
Section~\ref{sec:verify}.

Nothing else changes. Slots, array elements and pointer scaling all follow \code{W} and \code{W_SHIFT}, and
the frame keeps its 16-byte alignment, which the ILP32 ABI also requires. For example, the prologue of a
two-parameter function becomes:

\begin{center}
\begin{minipage}[t]{0.47\linewidth}
\begin{lstlisting}[language=riscv,numbers=none,title={\codetitle{RV64}}]
addi sp, sp, -32
sd ra, 24(sp)
sd s0, 16(sp)
addi s0, sp, 32
sd a0, -24(s0)
sd a1, -32(s0)
\end{lstlisting}
\end{minipage}\hfill
\begin{minipage}[t]{0.47\linewidth}
\begin{lstlisting}[language=riscv,numbers=none,title={\codetitle{RV32}}]
addi sp, sp, -16
sw ra, 12(sp)
sw s0, 8(sp)
addi s0, sp, 16
sw a0, -12(s0)
sw a1, -16(s0)
\end{lstlisting}
\end{minipage}
\end{center}

Because the whole HW3 suite can now run in both modes, the RV32 port is verified by the same 15 programs
as HW3 (\code{TARGET=rv32 make test}), on Spike with \code{--isa=rv32imac} and a 32-bit proxy kernel.

\section{The \texttt{uint32\_t} Type}

\subsection{Scanner and grammar}

\code{uint32_t} is a new token, and the grammar accepts it as a base type:

\snippet[language=flex]{../../scanner.l}{"uint32_t"}{+1}
\snippetafter{start program}
\snippet[language=yacc]{../../parser.y}{^base_type}{@rule}
\snippetafter{start program}
\snippet[language=yacc]{../../parser.y}{^sign}{@rule}

The semantic value of a type is a single flag: 1 for an unsigned type (\code{unsigned} \ldots or
\code{uint32_t}), 0 otherwise. A declaration copies the flag to every declarator, and \code{sym_add} stores it in
the symbol table:

\snippetafter{start program}
\snippet[language=yacc]{../../parser.y}{^local_decl}{@rule}
\snippet[language=yacc]{../../parser.y}{s->is_unsigned = where}{+1}

Function definitions, parameters and casts set the same flag, so return values and \code{(unsigned)x} are
covered as well.

\subsection{Which expressions are unsigned}

Addition, subtraction, multiplication and the bitwise operators give the same bits whether the operands
are signed or not. Only comparison, division, remainder and right shift differ, so the generator needs to
know the signedness of an expression just for those. \code{is_unsigned_expr} follows C's usual arithmetic
conversions:

\snippet[language=yacc]{../../parser.y}{static int is_unsigned_expr}{@end}

\begin{itemize}
  \item a variable, array element, or dereferenced variable is unsigned if it was declared so;
  \item a binary arithmetic or bitwise operation is unsigned if \emph{either} operand is;
  \item a comparison or logical operation yields a plain \code{int} (0 or 1), so it is signed;
  \item a shift takes the type of its \emph{left} operand only;
  \item every intrinsic returns \code{uint32_t}.
\end{itemize}

\subsection{Choosing the instruction}

\snippet[language=yacc]{../../parser.y}{int uns = }{const char *slt}

For comparisons and division the operation is unsigned if either side is; for \code{>>} only the left side
counts. The flag then selects the instruction:

\begin{center}
\begin{tabularx}{0.8\linewidth}{@{}lllL@{}}
\toprule
Operator & Signed & Unsigned & Example where they differ \\
\midrule
\code{< > <= >=} & \code{slt} & \code{sltu} & \code{0xFFFFFFFF > 1} \\
\code{/} & \code{div} & \code{divu} & \code{0xFFFFFFFF / 3} \\
\code{\%} & \code{rem} & \code{remu} & \code{0xFFFFFFFF \% 7} \\
\code{>>} & \code{sra} & \code{srl} & \code{0x80000000 >> 28} \\
\bottomrule
\end{tabularx}
\end{center}

\note{On RV64 a \code{uint32_t} is held in a 64-bit register, so wrap-around at $2^{32}$ is not modelled
(\code{0 - 1} becomes $2^{64} - 1$). Unsigned values are exact only with \code{--rv32}, which is the only
mode the Bonus targets.}

\section{Intrinsics}

\subsection{Recognising a call}

The intrinsics need no grammar changes: syntactically they are ordinary function calls. The generator keeps
a table that maps each name to its instruction:

\snippet[language=yacc]{../../parser.y}{static const Intrinsic INTRINSICS}{@end}

\code{gen_call} looks the name up first. A match goes to \code{gen_intrinsic} and emits a single instruction
instead of \code{jal}; any other name is an ordinary call (see the HW3 report).

\subsection{Emitting the instruction}

\snippet[language=yacc]{../../parser.y}{static void gen_intrinsic}{@end}

Evaluation follows the same stack-machine rule as a binary operator: the first argument is pushed, the
second is evaluated into \code{t0}, and the first is popped into \code{t1}. The instruction is then
\code{<op> t0, t1, t0}: \code{rd} is \code{t0} as for any expression, \code{rs1} is the first argument and
\code{rs2} the second. This order matters because \code{uksub8} and \code{ucmplt8} are not commutative. It
matches the TA's hand-written sample \code{ukadd8 a5, a4, a5}, where \code{a4} holds the first argument.
The generator also checks that exactly two arguments are given.

For \code{return __rv__ukadd8(a, b);} the RV32 output is:

\begin{lstlisting}[language=riscv,title={\codetitle{codegen.S (RV32), body of the function}}]
  addi t0, s0, -12          # a
  lw t0, 0(t0)
  addi sp, sp, -16          # push a
  sw t0, 0(sp)
  addi t0, s0, -16          # b
  lw t0, 0(t0)
  lw t1, 0(sp)              # pop a into t1
  addi sp, sp, 16
  ukadd8 t0, t1, t0         # rd = t0, rs1 = a, rs2 = b
  mv a0, t0
\end{lstlisting}

Nested calls such as \code{__rv__ukadd8(__rv__uksub8(a, b), __rv__cmpeq8(a, b))} need nothing extra: each
inner intrinsic leaves its result in \code{t0} like any other expression.

\section{Verification Without the Board}
\label{sec:verify}

The final check is to run the output on a Corvette board. Without one, a P-extension program cannot run
anywhere: upstream Spike has dropped the draft P-extension, and GNU \code{as} does not accept
\code{ukadd8}. \code{make test-bonus} (script \file{bonus/tests/run.sh}) therefore checks each test program in
three separate ways.

\subsection{Reference results}

A C reference implementation of the four instructions, written from the RISC-V P-extension specification:

\snippet{../tests/pext_ref.h}{static inline uint32_t pext_lanes}{@end}

The host gcc compiles the test program with each \code{__rv__*} name mapped to this implementation and runs it
on 12 input pairs. The output is the expected result. The inputs include the specification's examples and,
for each byte lane, the saturation and equality edge cases:

\snippet{../tests/main.c}{static const uint32_t inputs}{@end}

\subsection{Running the generated code}

\code{./codegen --rv32 --pext-emulate} generates code that calls \code{__pext_<op>(rs1, rs2)} in place of each
P-extension instruction, with the operands in the same order as the real instruction
(\code{a0} = \code{rs1}, \code{a1} = \code{rs2}). This is the \code{pext_emulate} branch of \code{gen_intrinsic}
shown above. The emulation functions call the same reference implementation:

\snippet{../tests/pext_emulate.c}{uint32_t __pext_ukadd8}{+4}

The program is linked with the same \file{main.c}, run on RV32 Spike, and its output must equal the
reference results exactly. Everything except the four instructions --- the RV32 frame, argument passing,
operand order, unsigned operations --- is executed for real.

\subsection{Checking the real output}

\code{./codegen --rv32} (without emulation) produces the \file{codegen.S} that goes onto the board. The script
checks it statically:

\snippet[language=bash]{../tests/run.sh}{# 3. Real output}{not valid RV32 assembly}

\begin{itemize}
  \item the number of P-extension instructions equals the number of emulated calls in step 2, so no
        intrinsic is lost or duplicated;
  \item no call to an \code{__rv__} function remains;
  \item with the P-extension instructions replaced by \code{nop}, the rest assembles as valid RV32 code.
\end{itemize}

\subsection{Test programs}

\begin{center}
\begin{tabularx}{\linewidth}{@{}lL@{}}
\toprule
Test & Content \\
\midrule
\file{0.c} & the TA's example \\
\file{1.c} & three intrinsics chained through local variables \\
\file{2.c} & non-commutative operations, to check operand order \\
\file{3.c} & nested intrinsic calls \\
\file{4.c} & unsigned comparison, division and right shift mixed with an intrinsic; \textbf{CRLF line endings}, like the TA's original \file{0.c} \\
\bottomrule
\end{tabularx}
\end{center}

\file{.gitattributes} sets \code{-text} on \file{4.c} so that git never converts its line endings. Without
it, git (configured with \code{core.autocrlf=input}) would store the file with LF endings and the CRLF case
would silently disappear; this is what happened to the TA's \file{0.c}.

\subsection{Testing the tests}

To make sure the checks can actually fail, I edited the generated assembly by hand and confirmed that each
change makes the suite fail:

\begin{center}
\begin{tabularx}{\linewidth}{@{}lL@{}}
\toprule
Mutation & Caught by \\
\midrule
swap the operands of an intrinsic & \file{1.c}, \file{2.c}, \file{3.c} \\
\code{sltu} $\rightarrow$ \code{slt} & \file{4.c} \\
\code{divu} $\rightarrow$ \code{div} & \file{4.c} \\
\code{srl} $\rightarrow$ \code{sra} & \file{4.c} \\
\bottomrule
\end{tabularx}
\end{center}

The first versions of two tests did not catch their mutation. \file{2.c} computed an expression that happened
to be symmetric, so swapping operands gave the same result; in \file{4.c} the difference between signed and
unsigned comparison was absorbed by a later division. Both were rewritten until every mutation was caught.

\section{Remaining Work}

\begin{itemize}
  \item \textbf{Hardware semantics.} The reference implementation follows the published P-extension
        specification; whether it matches the Andes hardware exactly can only be confirmed on the board.
  \item \textbf{Andes assembler.} The output uses the same syntax as the TA's sample (\code{ukadd8 a5, a4, a5}),
        but it has not yet been assembled with the Andes toolchain.
\end{itemize}

To run on a board:

\begin{enumerate}
  \item Generate the assembly with \code{./codegen --rv32 < testcase.c}.
  \item Install the Corvette board package as described in \file{bonus/spec/Arduino_Env_Guide_Corvette-F1_T1.pdf},
        and add \code{-mext-dsp} to \file{platform.txt}.
  \item Put \file{HW3_Bonus.ino} and the generated \file{codegen.S} in the same sketch folder and upload it.
  \item Open the Serial Monitor (no line ending, 9600 baud) and enter pairs such as \code{65280, 65283}.
        The sketch should print ``The results are equal.''
\end{enumerate}

\end{document}
